
\subsubsection{6.1 Efficiency Without Convergence}

The 10.2\% FLOP reduction derives from architectural iteration reduction (T=2 vs T=3), not from attention patterns converging to stable distributions. The biological intuition that temporal processing enables refinement toward task-relevant patterns is not supported by these experiments.

\subsubsection{6.2 Iteration Count Dominates K/V Caching}

Reducing temporal steps (10.2\% reduction) provides larger savings than K/V caching (6.8\% reduction). K/V projection constitutes a smaller fraction of total attention cost than initially estimated.

\subsubsection{6.3 Constant Scaling}

The 5.31\% reduction is sequence-length-independent because the T\_steps ratio cancels in the FLOP ratio. Adaptive T\_steps schedules that increase with sequence length may be needed to achieve scaling efficiency gains.

\subsubsection{6.4 Limitations}

\begin{enumerate}
\item \textbf{L1:} Convergence tested on randomly-initialized model. Trained models may exhibit different entropy dynamics, but this was not tested.
\item \textbf{L2:} CPU-only execution. Memory profiling data is not available.
\item \textbf{L3:} Temporal reasoning tasks (P3 prediction) were not tested.
\item \textbf{L4:} Single architecture (16-layer, 512-dim transformer) tested.
\item \textbf{L5:} Perplexity values (~265,000-270,000) reflect random initialization, not trained model quality. Relative comparisons remain valid.
\end{enumerate}

